\chapter*{TÓM LƯỢC}
\addcontentsline{toc}{chapter}{TÓM LƯỢC}

Quản lý tiến độ và chi phí trong các dự án logistics quy mô lớn đòi hỏi giải quyết bài toán Lập lịch Dự án Hạn chế Tài nguyên (RCPSP). Các phương pháp lập lịch định định (CPM/PERT) hoặc các bộ giải chính xác (CP-SAT) thường gặp giới hạn về bùng nổ tổ hợp (lớp bài toán NP-hard) khi áp dụng trên đồ thị mạng lưới quy mô lớn, dẫn đến thời gian tính toán tăng theo cấp số nhân.

Đồ án đề xuất hệ thống \textbf{GLPO (Graph-Based Logistics Project Optimizer)} kết hợp mô hình học máy và thuật toán tối ưu hóa tổ hợp. Thành phần cốt lõi là một luồng xử lý AI (AI Pipeline) ứng dụng mạng nơ-ron đồ thị dị thể (Heterogeneous Graph Transformer - HGT) và cơ chế học tự giám sát (Masked Autoencoder - MAE). Mô hình này thực hiện trích xuất đặc trưng cấu trúc (structural embedding) của đồ thị công việc, đóng vai trò như một hàm heuristic định hướng nhằm thu hẹp không gian tìm kiếm cho bộ giải Google OR-Tools (CP-SAT). Phương pháp lai (hybrid approach) này giúp hội tụ nghiệm tối ưu Pareto cục bộ trong thời gian giới hạn đối với bài toán RCPSP đa mục tiêu.

Về mặt kiến trúc phần mềm, hệ thống được phân rã thành các dịch vụ độc lập. Tầng \textbf{Backend} ứng dụng FastAPI (Python 3.12) và ORM bất đồng bộ (SQLAlchemy 2.0) để tương tác với cơ sở dữ liệu PostgreSQL 15. Tầng \textbf{Frontend} sử dụng React, TypeScript và Vite nhằm trực quan hóa cấu trúc đồ thị định hướng không chu trình (DAG) và biểu đồ Gantt.

Toàn bộ hệ thống được đóng gói (containerized) bằng Docker, đáp ứng các yêu cầu về tái lập môi trường thực nghiệm và khả năng triển khai liên tục, cung cấp cơ sở kỹ thuật cho việc ứng dụng Deep Graph Learning vào tối ưu hóa chuỗi cung ứng.